\chapter{Migrating an XNA or FNA Game to CNA}
\label{ch:migration-guide}

This chapter is the practical payoff of the preceding parts: a checklist for judging, before you
start, how much of an existing XNA or FNA game's own code and content will simply work on
CNA, and which specific gaps this book has already named in detail are the ones likely to
cost you real time.

\section{Step one: is this a C\# port, or a from-scratch rewrite?}

Be clear about which kind of migration you are actually doing. CNA preserves XNA's API
shape, not its language --- porting an existing C\# XNA/FNA game means rewriting its logic in
C\texttt{++} against the equivalent CNA classes, following the same conventions
Chapter~\ref{ch:design-philosophy} described for CNA's own internal porting work:
matching class and method names exactly, using \texttt{sharp-runtime}'s type aliases rather
than raw C\texttt{++} primitives, and translating C\# properties into the
\texttt{getXProperty()} / \texttt{setXProperty()} convention. This is mechanical translation
work for code that only ever touched the public XNA surface, and genuinely difficult work for
code that reached into .NET-specific behavior (reflection, real garbage collection timing,
LINQ) that \texttt{sharp-runtime} (Chapter~\ref{ch:sharp-runtime-overview}) deliberately does
not reproduce.

\begin{lstlisting}
// Worked example: a small, typical XNA Update()/Draw() pair, and its
// mechanical CNA translation -- the shape every line maps to directly.

// --- Original C# (XNA/FNA) ---
// protected override void Update(GameTime gameTime)
// {
//     float dt = (float)gameTime.ElapsedGameTime.TotalSeconds;
//     position += velocity * dt;
//     if (Keyboard.GetState().IsKeyDown(Keys.Space))
//         velocity.Y = -jumpSpeed;
//     base.Update(gameTime);
// }
//
// protected override void Draw(GameTime gameTime)
// {
//     GraphicsDevice.Clear(Color.CornflowerBlue);
//     spriteBatch.Begin();
//     spriteBatch.Draw(texture, position, Color.White);
//     spriteBatch.End();
//     base.Draw(gameTime);
// }

// --- Mechanical CNA C++23 translation ---
void MyGame::Update(GameTime& gameTime)
{
    float dt = static_cast<float>(gameTime.getElapsedGameTimeProperty().getTotalSecondsProperty());
    position_ += velocity_ * dt;
    if (Keyboard::GetState().IsKeyDown(Keys::Space))
        velocity_.Y = -jumpSpeed_;
    Game::Update(gameTime);
}

void MyGame::Draw(const GameTime& gameTime)
{
    getGraphicsDeviceProperty().Clear(Color::CornflowerBlue);
    spriteBatch_->Begin();
    spriteBatch_->Draw(texture_, position_, Color::White);
    spriteBatch_->End();
    Game::Draw(gameTime);
}
\end{lstlisting}

\section{Step two: audit content against the two biggest named gaps}

Before touching gameplay code, check the game's own content against the two gaps this book
has returned to most often:

\begin{itemize}
  \item \textbf{Compiled \texttt{.fx} shader bytecode} (Chapter~\ref{ch:shader-fx-gap}).
        Any custom shader compiled into the original game's content pipeline will hit
        \cnaclass{Effect}'s bytecode constructor, which always throws. A game built entirely
        on the five stock effects and \cnaclass{SpriteBatch} (Chapter~\ref{ch:stock-effects})
        is unaffected; \texttt{cna-samples}' own 23-of-86 failure count
        (Chapter~\ref{ch:samples-and-examples}) is a real, checkable proxy for how common
        this gap actually is across real XNA content.
  \item \textbf{The \texttt{.xnb} content pipeline} (Chapter~\ref{ch:content-and-assets}).
        Real \texttt{.xnb} compiled assets do not load directly except through the specific
        readers that exist for them --- \cnaclass{Model} has a genuine, wire-compatible
        \texttt{.xnb} reader today (Chapter~\ref{ch:models-meshes}); most other asset types do
        not yet, and need converting to CNA's raw-file-plus-JSON-descriptor convention
        instead.
\end{itemize}

\section{Step three: check vertex layouts against the five recognized strides}

If the game does any custom 3D rendering, Chapter~\ref{ch:textures-rendertargets}'s vertex
layout finding is worth checking directly: every hardware renderer selects its GPU
vertex-attribute layout by raw byte stride, not by the \cnaclass{VertexDeclaration} you
actually wrote, and only five specific strides (16, 20, 24, 32, and 52 bytes, corresponding to
the standard \cnaclass{VertexPositionColor} / \texttt{Texture} / \texttt{ColorTexture} /
\texttt{NormalTexture} structs and one skinned layout) are guaranteed to render correctly
everywhere. A custom vertex format outside that set will degrade differently on each renderer
--- Chapter~\ref{ch:textures-rendertargets} gives the specific per-renderer fallback behavior
--- so this is worth checking before porting any custom mesh format, not after seeing
incorrect rendering on one specific renderer.

\begin{lstlisting}
// Worked example: computing a custom vertex format's stride by hand
// before porting it, the check this step actually means in practice.
struct MyCustomVertex
{
    Vector3 Position;  // 12 bytes
    Vector3 Normal;    // 12 bytes
    Vector2 TexCoord;  // 8 bytes
    Vector4 Tangent;   // 16 bytes
};
// sizeof(MyCustomVertex) == 48 bytes -- NOT one of the five recognized
// strides (16/20/24/32/52), so this format is not guaranteed to render
// identically across every renderer. Decide up front: drop the Tangent
// field to reach the 32-byte NormalTexture stride, or accept and test
// the per-renderer fallback behavior Chapter 17 documents for an
// unrecognized stride, rather than discovering the mismatch as a visual
// bug later.
static_assert(sizeof(MyCustomVertex) == 48);
\end{lstlisting}

\section[Step four: networking and gamer services needs]{Step four: check what your game actually needs from networking and gamer services}

If the original game used Xbox Live features, Chapter~\ref{ch:networking} and
Chapter~\ref{ch:gamerservices} give you the exact real-versus-stub breakdown: system-link
multiplayer, local achievements, and local leaderboards are genuinely functional; matchmaking,
Xbox Live friends, and most of the \cnaclass{Guide} system's UI are documented stubs with no
local equivalent to fall back to, because there is no local substitute for a service that
required a live Xbox Live connection to begin with. A game whose multiplayer was always
system-link/LAN-based (as \emph{Speedy Blupi} and \emph{Planet Blupi} were, per
Chapter~\ref{ch:blupi-case-study}) ports its networking with comparatively little friction;
a game built around Xbox Live matchmaking or friends lists does not have a real target to
port that part to at all. One specific, easy-to-miss gotcha worth checking before porting any
achievement-toast UI: \cnaclass{Achievement} objects \cnaclass{SignedInGamer::GetAchievements()}
returns carry real \texttt{Key} / \texttt{IsEarned} / \texttt{EarnedDateTime} data, but
\texttt{Name}, \texttt{Description}, and \texttt{GamerScore} stay at fixed, empty defaults ---
the original game's own achievement catalog metadata never had a local source of truth to begin
with (Chapter~\ref{ch:gamerservices}'s own worked example covers this precisely). A game whose
achievement-earned toast reads directly from the returned \cnaclass{Achievement} object's own
\texttt{Name} / \texttt{GamerScore} fields will port and compile cleanly, then silently render a
blank name and \texttt{+0G} for every single achievement --- the fix is maintaining your own
small achievement-metadata table keyed by \texttt{Key}, not something this namespace can supply
for you.

A second, equally easy-to-miss gotcha in the same namespace, worth checking before porting any
code that relies on a specific session capacity: every \cnaclass{NetworkSession::Create}/%
\texttt{BeginCreate} overload's own \texttt{maxGamers} argument is silently never honored ---
the constructed session's real \texttt{MaxGamers} is always a hardcoded \texttt{69}, matching
genuine historical Xbox~360/Windows XNA~4.0 behavior for this exact call shape, not a CNA bug
(Chapter~\ref{ch:networking}'s own worked example covers the precise mechanism). A game that
calls \texttt{Create(SystemLink, maxLocalGamers, 8)} expecting an eight-player cap and enforces
nothing further at the application level will accept a ninth, tenth, or sixty-ninth join before
the framework itself ever refuses one --- if the original game's own design assumed the
requested cap was load-bearing, that enforcement needs to move into the game's own code
explicitly, not stay implicit in the constructor call.

\section{Step five: pick and verify your graphics renderer}

Part~IV gave a full per-renderer account; for a porting decision specifically,
three questions matter most. Is the game 2D-only? \texttt{SDL\_RENDERER} or \texttt{OPENGLES3}
are both mature, exhaustively pixel-verified choices (Chapters~\ref{ch:sdl-renderer}
and~\ref{ch:easygl-backend}). Does it need real 3D with stock effects only?
\texttt{OPENGLES3} or \texttt{VULKAN} are the most broadly verified choices, with BGFX close
behind (Chapters~\ref{ch:easygl-backend}--\ref{ch:bgfx-backend}). Does correctness against
real, historical XNA rendering matter more than portability --- for instance, a game whose
original visual reference is what a screenshot-comparison test suite checks against?
the \texttt{DIRECTX9} renderer's oracle-verified Direct3D~9 pixel authenticity
(Chapter~\ref{ch:direct3d-backends}) is the only renderer built specifically for that bar,
at the cost of being Windows-only
(cross-compiled and Wine-verified, per Chapter~\ref{ch:windows-wine}).

\section{Step six: use xna4-spec, not memory, to check what you ported}

Once a subsystem is ported, Chapter~\ref{ch:xna4-spec-auditing}'s auditing method --- checking
against the real XNA~4.0 specification directly, not against FNA's own source or against
recollection --- is the single most reliable way to confirm you did not miss a member or
introduce an unmarked extension, for exactly the reason Chapter~\ref{ch:media-system}'s
\cnaclass{Song} finding illustrated: FNA itself is not always a perfectly faithful stand-in
for the real spec.

\section{Step seven: troubleshooting the build and the first run}

The six steps above assume the port compiles and the graphics renderer actually starts.
\texttt{docs/migration-guide.md} in the \texttt{cna} repository (Task 486, the source this
whole chapter draws its structure from) closes with two practical troubleshooting tables ---
real configure-time and runtime failures a porter is likely to hit, each checked directly
against the current \texttt{CMakeLists.txt}/source rather than assumed from general CMake or
XNA knowledge. Both are reproduced here, condensed, since they are exactly the kind of
concrete, checkable-cost information this chapter's own closing section already asks you to
value over a general sense of the project's maturity.

\subsection{Configure-time failures}

\begin{center}
\begin{longtable}{>{\raggedright\arraybackslash}p{0.32\textwidth}>{\raggedright\arraybackslash}p{0.58\textwidth}}
\toprule
\textbf{Symptom} & \textbf{Cause and fix} \\
\midrule
\endhead
\texttt{CNA: Missing sibling repository 'sharp-runtime'} / \texttt{'easy-gl'} & These are
  sibling checkouts next to \texttt{cna}, \emph{not} git submodules ---
  \texttt{git submodule update --init} will not fetch them. Clone the missing repo as a
  sibling directory instead; the real error message prints the exact URL. \\[4pt]
\texttt{Missing vendored 'SDL' \dots} right after a fresh clone & Downloaded a ZIP/release
  archive instead of cloning with Git, or forgot to init submodules --- vendored
  \texttt{third\_party/SDL} / \texttt{SDL\_image} / \texttt{SDL\_mixer} are empty. Run
  \texttt{git submodule update --init --recursive}, or configure with
  \texttt{-DCNA\_USE\_SYSTEM\_SDL=ON} to use system SDL3 packages instead. \\[4pt]
\texttt{Could not find a package configuration file for Vulkan} (\texttt{VULKAN} renderer) &
  No Vulkan SDK/loader installed. Install your distribution's \texttt{vulkan-sdk}/%
  \texttt{libvulkan-dev} package, or the LunarG SDK, before configuring. \\[4pt]
CMake \texttt{FetchContent} hangs or fails cloning \texttt{bgfx.cmake} (\texttt{BGFX} renderer)
  & \texttt{BGFX} fetches \texttt{bgfx.cmake} from GitHub at configure time --- needs network
  access, and the first fetch can be slow. Retry, or pre-seed a local clone and point
  \texttt{FETCHCONTENT\_SOURCE\_DIR\_BGFX\_CMAKE} at it. \\[4pt]
\cnaclass{VideoPlayer}-related link errors & Missing FFmpeg dev packages --- hits every
  renderer, not just video code. Install \texttt{libavcodec-dev libavformat-dev
  libavutil-dev libswresample-dev}. \\
\bottomrule
\end{longtable}
\end{center}

\subsection{Runtime failures}

\begin{center}
\begin{longtable}{>{\raggedright\arraybackslash}p{0.32\textwidth}>{\raggedright\arraybackslash}p{0.58\textwidth}}
\toprule
\textbf{Symptom} & \textbf{Cause and fix} \\
\midrule
\endhead
Any 3D call throws \texttt{std::runtime\_error} & You are on \texttt{SDL\_RENDERER} --- not a
  bug, it is 2D-only by design (Chapter~\ref{ch:sdl-renderer}). Switch to \texttt{OPENGLES3},
  \texttt{VULKAN}, or \texttt{BGFX} for 3D. \\[4pt]
\texttt{bgfx::init failed}, or a native-window-handle error on \texttt{BGFX} & No usable
  native window handle for \texttt{bgfx}'s chosen renderer under the current SDL video
  driver. Check \texttt{SDL\_GetCurrentVideoDriver()}'s output (printed in the real error
  message); try forcing a renderer explicitly (see below). \\[4pt]
\texttt{BGFX} picks an unavailable graphics API (e.g.\ Vulkan on a machine with no Vulkan
  ICD) & \texttt{bgfx}'s own renderer auto-detection picked something unavailable. Set the
  \texttt{CNA\_BGFX\_RENDERER} environment variable to force one explicitly: \texttt{AUTO},
  \texttt{OPENGL}, \texttt{OPENGLES}, \texttt{VULKAN}, \texttt{METAL}, \texttt{DIRECT3D11},
  \texttt{DIRECT3D12}, or \texttt{NOOP} (case-insensitive; an unsupported value throws
  immediately with this exact list). \\[4pt]
Blank/black window, or an OpenGL-based renderer fails to create a context in CI/headless
  environments & No real X server or GPU driver --- common under \texttt{Xvfb} with only a
  software (\texttt{llvmpipe}) GL driver (Chapters~\ref{ch:sdl-renderer}
  and~\ref{ch:easygl-backend} both now include a real, confirmed-working example of exactly
  this setup). Expected in headless CI; tests still run against software rendering. If an
  entire run fails rather than just individual pixel-tolerance edge cases, confirm a display
  is actually reachable (the \texttt{DISPLAY} environment variable, or
  \texttt{-DCNA\_TEST\_DISPLAY=:99} at configure time to point the whole \texttt{ctest} suite
  at a specific \texttt{Xvfb} display). \\[4pt]
A pixel-comparison example or test fails right after an intentional rendering change &
  Golden images are stale references, not a bug. Re-run the specific test with
  \texttt{CNA\_UPDATE\_GOLDEN=1} set to regenerate the golden PNG, then review the diff before
  committing it (Chapter~\ref{ch:verification-methodology} covers this golden-image mechanism
  in full). \\[4pt]
Flaky test failures that do not reproduce in isolation & \textbf{Never run multiple
  renderers' full \texttt{ctest} suites concurrently} --- confirmed to cause transient
  GPU/driver-contention false failures. Run renderer test suites sequentially, and re-run any
  single anomalous test in isolation (\texttt{ctest -R <TestName>}) before treating it as a
  real regression, not a race. \\
\bottomrule
\end{longtable}
\end{center}

Before filing anything as a new CNA bug during a port, check three sources in this order:
\texttt{docs/xna-4-api-coverage.md}'s per-renderer support table and its ``Known deviations
from XNA/FNA'' list, then \texttt{NEXT.md}'s own actively-maintained current bug section ---
most currently-open gaps already have a task number and a documented root cause. If a symptom
genuinely is not covered by any of the three, it is more likely a real, new finding worth its
own report than a known, already-tracked limitation.

\section{A realistic expectation, stated plainly}

Every chapter in this book that has named a real, still-open gap has also named its measured
cost where one exists --- 23 of 86 samples, a specific 94-test blast radius, a specific
still-open leak in \texttt{NetworkSession::BeginCreate}. Use those numbers, not a general
sense of the project's maturity, to judge your own porting timeline: a 2D, stock-effects-only,
system-link-multiplayer game is a comparatively mechanical port today; a game leaning on
custom compiled shaders, Xbox Live matchmaking, or an unusual vertex format will hit named,
real, and in some cases still-unresolved gaps, and this book has tried throughout to tell you
exactly which ones, rather than leave you to discover them yourself.
